\documentclass[12pt]{article}
\usepackage[utf8]{inputenc}
\usepackage[T1]{fontenc}
\usepackage[margin=1in]{geometry}
\usepackage{setspace}
\usepackage{xcolor}
\usepackage[colorlinks=true, urlcolor=blue!60!black, linkcolor=blue!60!black]{hyperref}
\setlength{\parindent}{2em}
\setlength{\parskip}{1.25em}
\renewcommand{\baselinestretch}{1.0}
\begin{document}
\begin{center}
{\LARGE \textit{Grutter v. Bollinger}, 539 U.S. 306 (2003)\par}
\vspace{0.8em}
{\large Case Brief\par}
\end{center}
\vspace{1em}

\section{Memory Jogger}
In \href{https://www.courtlistener.com/opinion/130155/grutter-v-bollinger/}{\textit{Grutter v. Bollinger}}, a 5--4 Supreme Court upheld the University of Michigan Law School's holistic, race-conscious admissions policy under strict scrutiny, holding that attaining the educational benefits of student body diversity is a compelling state interest---a framework later effectively ended twenty years later in \href{https://supreme.justia.com/cases/federal/us/600/20-1199/}{\textit{Students for Fair Admissions, Inc. v. Harvard}}, 600 U.S. 181 (2023).

\section{Detailed Case Facts}
The University of Michigan Law School, one of the nation's top public law schools, received more than 3,500 applications each year for an entering class of roughly 350 students (\href{https://supreme.justia.com/cases/federal/us/539/306/}{Justia Opinion}). In 1992, a faculty committee chaired by Professor Richard Lempert drafted a written admissions policy, unanimously adopted by the faculty, designed to comply with Justice Lewis Powell's opinion in \textit{Regents of the University of California v. Bakke}, 438 U.S. 265 (1978) (\href{https://www.law.cornell.edu/supremecourt/text/02-241}{Cornell LII}).

The 1992 policy required admissions officials to evaluate each applicant holistically based on the entire file---undergraduate GPA, LSAT score, personal statement, letters of recommendation, an essay on how the applicant would contribute to Law School life and diversity, and ``soft variables'' such as recommenders' enthusiasm, undergraduate institution quality, and course difficulty (\href{https://supreme.justia.com/cases/federal/us/539/306/}{539 U.S. at 315}). While recognizing many bases for diversity, the policy reaffirmed a special commitment to enrolling a ``critical mass'' of historically underrepresented minority students---specifically African-American, Hispanic, and Native-American applicants---so that minority students would not feel isolated or like spokespersons for their race and so classroom discussion would benefit from diverse perspectives (\href{https://supreme.justia.com/cases/federal/us/539/306/}{539 U.S. at 316, 318--319}). Admissions Director Dennis Shields testified that he consulted ``daily reports'' tracking the class's racial composition during the admissions season, though he denied seeking any fixed number or percentage (\href{https://supreme.justia.com/cases/federal/us/539/306/}{539 U.S. at 318}).

In 1996, Barbara Grutter, a white Michigan resident with a 3.8 undergraduate GPA and a 161 LSAT score, applied to the Law School, was waitlisted, and was subsequently rejected (\href{https://supreme.justia.com/cases/federal/us/539/306/}{539 U.S. at 316}). In December 1997, Grutter sued the Law School, the Regents of the University of Michigan, Lee Bollinger (former Law School Dean and then University President), Dean Jeffrey Lehman, and Admissions Director Dennis Shields, alleging racial discrimination in violation of the Equal Protection Clause of the Fourteenth Amendment, Title VI of the Civil Rights Act of 1964 (42 U.S.C. \S~2000d), and 42 U.S.C. \S~1981 (\href{https://supreme.justia.com/cases/federal/us/539/306/}{539 U.S. at 316--317}). At the 15-day bench trial, Grutter's expert Dr. Kinley Larntz testified that minority status was ``an extremely strong factor'' (though not the predominant factor) in admissions, while the Law School's expert Dr. Stephen Raudenbush testified that a race-blind system in 2000 would have cut underrepresented minority admission rates from 35\% to 10\% and reduced underrepresented minority enrollment from 14.5\% to 4\% of the entering class (\href{https://supreme.justia.com/cases/federal/us/539/306/}{539 U.S. at 320}).

\section{Procedural History}
\begin{itemize}
    \item \textbf{U.S. District Court for the Eastern District of Michigan (Mar. 27, 2001):} Judge Bernard A. Friedman certified the plaintiff class, conducted a 15-day bench trial, held that student body diversity was not a compelling state interest and that the Law School's ``critical mass'' policy was not narrowly tailored (deeming it practically indistinguishable from a quota), and enjoined the Law School from using race in admissions, \textit{Grutter v. Bollinger}, 137 F. Supp. 2d 821 (E.D. Mich. 2001) (\href{https://supreme.justia.com/cases/federal/us/539/306/}{Justia Opinion}; \href{https://en.wikipedia.org/wiki/Grutter_v._Bollinger}{Wikipedia}).
    \item \textbf{U.S. Court of Appeals for the Sixth Circuit (May 14, 2002):} After staying the injunction pending appeal, the Sixth Circuit sitting en banc reversed 5--4 and vacated the injunction, holding that Justice Powell's \textit{Bakke} opinion was binding under \textit{Marks v. United States}, 430 U.S. 188 (1977), and that the policy was narrowly tailored, 288 F.3d 732 (6th Cir. 2002) (\href{https://supreme.justia.com/cases/federal/us/539/306/}{539 U.S. at 321}).
    \item \textbf{Supreme Court of the United States (June 23, 2003):} Granted certiorari to resolve a circuit split between the Fifth and Ninth Circuits, 537 U.S. 1043 (2002); heard oral argument on April 1, 2003 (alongside companion case \textit{Gratz v. Bollinger}, 539 U.S. 244 (2003), which struck down Michigan's point-based undergraduate admissions policy); and affirmed the Sixth Circuit in \href{https://www.courtlistener.com/opinion/130155/grutter-v-bollinger/}{539 U.S. 306 (2003)}. Rehearing was denied at 539 U.S. 982 (2003) (\href{https://en.wikipedia.org/wiki/Grutter_v._Bollinger}{Wikipedia}).
\end{itemize}

\section{Judicial Votes}
The Supreme Court affirmed the Sixth Circuit by a \textbf{5--4} vote (\href{https://supreme.justia.com/cases/federal/us/539/306/}{Justia Opinion}; \href{https://www.law.cornell.edu/supremecourt/text/02-241}{Cornell LII}):
\begin{itemize}
    \item \textbf{Opinion of the Court (5 Justices joining in full):} Justice Sandra Day O'Connor (author), joined by Justices John Paul Stevens, David H. Souter, Ruth Bader Ginsburg, and Stephen G. Breyer.
    \begin{itemize}
        \item \textit{Partial Joinder:} Justices Antonin Scalia and Clarence Thomas joined the Court's opinion \textit{in part}, solely insofar as it is consistent with Part VII of Justice Thomas's opinion (confirming that racial balancing or discrimination among preferred minority groups remains unlawful, and agreeing with the 25-year durational limit in the final sentence of Part III).
    \end{itemize}
    \item \textbf{Concurring Opinion:} Justice Ruth Bader Ginsburg (author), joined by Justice Stephen G. Breyer.
    \item \textbf{Concurrences in Part and Dissents in Part (2 Justices):}
    \begin{itemize}
        \item Justice Antonin Scalia (author), joined by Justice Clarence Thomas (also joining Rehnquist's dissent and Parts I--VII of Thomas's opinion).
        \item Justice Clarence Thomas (author), joined by Justice Antonin Scalia as to Parts I--VII (also joining Rehnquist's dissent and Scalia's opinion).
    \end{itemize}
    \item \textbf{Dissents (4 Justices total voting to reverse):}
    \begin{itemize}
        \item Chief Justice William H. Rehnquist (author), joined by Justices Antonin Scalia, Anthony M. Kennedy, and Clarence Thomas.
        \item Justice Anthony M. Kennedy (author, dissenting separately while also joining Rehnquist's dissent in full).
    \end{itemize}
\end{itemize}

\section{Holding}
\textbf{Judgment:} Affirmed. The University of Michigan Law School's race-conscious admissions policy did not violate the Equal Protection Clause of the Fourteenth Amendment, Title VI of the Civil Rights Act of 1964, or 42 U.S.C. \S~1981 (\href{https://www.courtlistener.com/opinion/130155/grutter-v-bollinger/}{CourtListener}; \href{https://supreme.justia.com/cases/federal/us/539/306/}{539 U.S. at 343--344}).

\textbf{Rule Announced:} Although all governmental racial classifications trigger strict scrutiny, strict scrutiny is not ``strict in theory, but fatal in fact'' (\href{https://supreme.justia.com/cases/federal/us/539/306/}{539 U.S. at 326}). In public higher education, attaining the educational benefits that flow from a diverse student body is a compelling state interest. A race-conscious admissions program is narrowly tailored to that interest if it: (1) uses race flexibly as a ``plus'' factor within a highly individualized, holistic review rather than a quota or mechanical point system; (2) seriously considers workable race-neutral alternatives in good faith without requiring sacrifice of academic selectivity; (3) does not unduly harm nonminority applicants; and (4) is limited in duration (\href{https://supreme.justia.com/cases/federal/us/539/306/}{539 U.S. at 333--343}).

\textbf{Holding vs. Dicta:}
\begin{itemize}
    \item \textbf{Holding:} (1) Student body diversity is a compelling governmental interest in higher education; (2) the Law School's holistic, individualized admissions program satisfied narrow tailoring under strict scrutiny; and (3) because the policy satisfied the Equal Protection Clause, Grutter's coextensive Title VI and \S~1981 claims also failed (\href{https://supreme.justia.com/cases/federal/us/539/306/}{539 U.S. at 325, 334, 343--344}).
    \item \textbf{Dicta:} Justice O'Connor's famous statement---``We expect that 25 years from now, the use of racial preferences will no longer be necessary to further the interest approved today'' (\href{https://supreme.justia.com/cases/federal/us/539/306/}{539 U.S. at 343})---was framed as an expectation (though Justices Thomas and Scalia treated it as a binding 25-year sunset holding, \href{https://supreme.justia.com/cases/federal/us/539/306/}{539 U.S. at 351, 375}). The Court also declined to decide whether Justice Powell's \textit{Bakke} opinion was binding under \textit{Marks}, rendering discussion of \textit{Marks} non-essential (\href{https://supreme.justia.com/cases/federal/us/539/306/}{539 U.S. at 325}).
\end{itemize}

\section{Analysis of Opinions}

\subsection{Majority Opinion (Justice O'Connor)}
Justice O'Connor began by endorsing Justice Powell's view in \textit{Bakke} directly, bypassing the lower-court tangle over \textit{Marks} (\href{https://supreme.justia.com/cases/federal/us/539/306/}{539 U.S. at 325}). Reaffirming \textit{Adarand Constructors, Inc. v. Pe\~{n}a}, 515 U.S. 200 (1995), she emphasized that all racial classifications face strict scrutiny, but ``[c]ontext matters'' (\href{https://supreme.justia.com/cases/federal/us/539/306/}{539 U.S. at 326--327}). Grounded in First Amendment academic freedom traditions, the Court deferred to the Law School's educational judgment that diversity is essential to its mission (\href{https://supreme.justia.com/cases/federal/us/539/306/}{539 U.S. at 328--329}). Drawing heavily on amicus briefs from Fortune 500 businesses (such as 3M and General Motors) and retired military leaders, O'Connor found the benefits of diversity ``not theoretical but real'': promoting cross-racial understanding, dismantling stereotypes, preparing graduates for a global workforce, ensuring a diverse military officer corps, and keeping the path to national leadership ``visibly open'' to all races (\href{https://supreme.justia.com/cases/federal/us/539/306/}{539 U.S. at 330--333}).

On narrow tailoring, O'Connor distinguished Michigan's holistic law school review from the mechanical 20-point undergraduate bonus struck down the same day in \textit{Gratz v. Bollinger}, 539 U.S. 244 (2003) (\href{https://supreme.justia.com/cases/federal/us/539/306/}{539 U.S. at 337}). Seeking a ``critical mass'' was a permissible flexible goal rather than an unconstitutional quota, evidenced by minority enrollment fluctuating between 13.5\% and 20.1\% from 1993 to 1998 and by the frequent admission of nonminority applicants with lower grades and test scores than rejected minority applicants (\href{https://supreme.justia.com/cases/federal/us/539/306/}{539 U.S. at 335--338}). Narrow tailoring did not require exhausting race-neutral alternatives like lotteries, lowered standards, or undergraduate ``percentage plans'' that would sacrifice academic selectivity or holistic review (\href{https://supreme.justia.com/cases/federal/us/539/306/}{539 U.S. at 339--340}). Finally, O'Connor insisted that race-conscious policies must have a ``logical end point'' and expressed the Court's expectation that racial preferences would no longer be necessary in 25 years (\href{https://supreme.justia.com/cases/federal/us/539/306/}{539 U.S. at 342--343}).

\subsection{Concurring Opinion (Justice Ginsburg, joined by Justice Breyer)}
Justice Ginsburg agreed that affirmative action must be temporally limited, citing international human rights treaties (ICERD and CEDAW), but cautioned against treating 25 years as a firm expiration date (\href{https://supreme.justia.com/cases/federal/us/539/306/}{539 U.S. at 344--346}). Pointing to persistent racial bias and stark K--12 school segregation and funding disparities (noting that in 2000--2001, 71.6\% of Black children and 76.3\% of Hispanic children attended majority-minority schools), Ginsburg concluded that ``one may hope, but not firmly forecast'' that equal opportunity would allow sunsetting affirmative action within a single generation (\href{https://supreme.justia.com/cases/federal/us/539/306/}{539 U.S. at 345--346}).

\subsection{Concurrence in Part and Dissent in Part (Justice Scalia, joined by Justice Thomas)}
Joining Chief Justice Rehnquist's dissent and Parts I--VII of Justice Thomas's opinion, Justice Scalia ridiculed the ``mystical `critical mass' justification'' as a sham masking racially proportionate admissions (\href{https://supreme.justia.com/cases/federal/us/539/306/}{539 U.S. at 346--347}). He mocked ``cross-racial understanding'' as a generic ``lesson of life rather than law'' (neither graded on law transcripts nor tested on the bar exam) that could just as easily justify racial discrimination in civil service or private employment (\href{https://supreme.justia.com/cases/federal/us/539/306/}{539 U.S. at 347--348}). Scalia further warned that the ``\textit{Grutter}--\textit{Gratz} split double header'' was perversely designed to spawn endless future litigation over de facto quotas, holistic review, and campus self-segregation (\href{https://supreme.justia.com/cases/federal/us/539/306/}{539 U.S. at 348--349}).

\subsection{Concurrence in Part and Dissent in Part (Justice Thomas, joined by Justice Scalia as to Parts I--VII)}
Opening with Frederick Douglass's 1865 plea---``Do nothing with us!''---Justice Thomas argued that the Constitution is strictly colorblind and ``abhors classifications based on race'' (\href{https://supreme.justia.com/cases/federal/us/539/306/}{539 U.S. at 349--350, 353, 378}). He contended that only ``pressing public necessity'' (such as national security or remedying an institution's own past discrimination) can satisfy strict scrutiny, and that Michigan had no compelling interest in maintaining a public law school at all---much less an ``elite'' one whose graduates largely leave the state---or in achieving a classroom ``aesthetic'' (\href{https://supreme.justia.com/cases/federal/us/539/306/}{539 U.S. at 351--361}). Thomas attacked the majority's deference to university administrators as antithetical to strict scrutiny, pointed to California's Proposition 209 experience at UC Berkeley's Boalt Hall as proof that race-neutral admissions work, and argued that preferences harm their intended beneficiaries through academic ``mismatch'' and racial stigma (\href{https://supreme.justia.com/cases/federal/us/539/306/}{539 U.S. at 362--373}). In Part VII, Thomas concurred solely to underscore that discrimination among minority groups remains unlawful and that the Court's 25-year limit meant racial preferences would be illegal in 2028 (\href{https://supreme.justia.com/cases/federal/us/539/306/}{539 U.S. at 374--378}).

\subsection{Dissenting Opinion (Chief Justice Rehnquist, joined by Justices Scalia, Kennedy, and Thomas)}
Chief Justice Rehnquist focused on narrow tailoring, arguing that ``[s]tripped of its `critical mass' veil, the Law School's program is revealed as a naked effort to achieve racial balancing'' (\href{https://supreme.justia.com/cases/federal/us/539/306/}{539 U.S. at 379}). Reviewing 1995--2000 admissions data, Rehnquist highlighted that the Law School admitted 91--108 African-American students, 47--56 Hispanic students, and only 13--19 Native-American students each year without explaining why ``critical mass'' required vastly different numbers across those three groups (\href{https://supreme.justia.com/cases/federal/us/539/306/}{539 U.S. at 381--383}). Presenting three statistical tables, he demonstrated a tight year-by-year correlation between each group's share of the applicant pool and its share of admitted students (e.g., African-Americans were 9.7\% of applicants and 9.4\% of admittees in 1995; 7.5\% of applicants and 7.3\% of admittees in 2000), proving forbidden racial balancing (\href{https://supreme.justia.com/cases/federal/us/539/306/}{539 U.S. at 383--386}). He also faulted the policy for lacking any concrete durational limit (\href{https://supreme.justia.com/cases/federal/us/539/306/}{539 U.S. at 386--387}).

\subsection{Dissenting Opinion (Justice Kennedy)}
Justice Kennedy accepted Justice Powell's \textit{Bakke} premise that student body diversity can be a compelling interest, but dissented because the majority abandoned genuine strict scrutiny in evaluating how Michigan implemented that goal (\href{https://supreme.justia.com/cases/federal/us/539/306/}{539 U.S. at 387--388}). Pointing to the narrow 0.3\% fluctuation band in enrolled minority percentages from 1995 to 1998 (13.5\% to 13.8\%) and admissions officers' heavy reliance on ``daily reports'' late in the admissions cycle, Kennedy argued that the pursuit of critical mass subverted individualized review and made race outcome-determinative for the final 15--20\% of seats (\href{https://supreme.justia.com/cases/federal/us/539/306/}{539 U.S. at 389--393}). ``Deference is antithetical to strict scrutiny, not consistent with it,'' Kennedy concluded (\href{https://supreme.justia.com/cases/federal/us/539/306/}{539 U.S. at 394}).

\section{Examples: Future Applications}
\begin{enumerate}
    \item \textbf{Same-Side (Valid under \textit{Grutter}): Holistic Public Medical School Admissions.} A state university medical school evaluates every applicant individually across MCAT scores, GPA, clinical service, personal hardship essays, and demographic background, treating race solely as one non-mechanical ``plus'' factor among many without quotas or fixed numerical targets, after documenting that race-neutral socioeconomic preferences failed to produce meaningful diversity. Under \textit{Grutter} (and later \textit{Fisher v. University of Texas at Austin} (\textit{Fisher II}), 579 U.S. 365 (2016)), this policy would be upheld---though it would now be invalid following \href{https://supreme.justia.com/cases/federal/us/600/20-1199/}{\textit{Students for Fair Admissions v. Harvard}} (2023).
    \item \textbf{Same-Side (Valid under \textit{Grutter}): Race-Neutral ``Top Ten Percent'' Plan Supplemented by Holistic Review.} A flagship public university admits the top 10\% of every in-state high school class automatically and fills the remaining 25\% of seats through a holistic file review where race is one contextual sub-factor within a broad personal-achievement index, coupled with periodic five-year reviews to assess whether race-conscious review remains necessary. Under \textit{Grutter}, this satisfies both compelling interest and narrow tailoring.
    \item \textbf{Opposite-Side (Invalid even under \textit{Grutter}): Mechanical Point-Bonus or Separate-Track Admissions.} A public college automatically adds 15 points on a 100-point admissions index to every applicant from an underrepresented racial minority group, or routes minority applicants to a separate admissions committee with reserved seats. Just like \textit{Gratz v. Bollinger}, 539 U.S. 244 (2003), and \textit{Bakke}, this fails \textit{Grutter}'s narrow-tailoring mandate because it lacks individualized consideration and insulates applicants from unified competition (\href{https://supreme.justia.com/cases/federal/us/539/306/}{539 U.S. at 334, 337}).
    \item \textbf{Opposite-Side (Invalid even under \textit{Grutter}): Municipal Police Promotion Preferences Justified by ``Workplace Diversity.''} A city police department gives a racial plus factor to minority sergeants seeking promotion to lieutenant solely to achieve a ``critical mass'' of diverse leadership and foster ``cross-racial understanding'' in the precinct, without any finding of prior departmental discrimination. Because \textit{Grutter}'s diversity deference was tied specifically to universities' First Amendment academic-freedom niche (\href{https://supreme.justia.com/cases/federal/us/539/306/}{539 U.S. at 328--329}), courts refuse to extend it to non-educational employment preferences.
    \item \textbf{Fence-Sitter (Close Case): Race-Neutral Proxy Policies Designed with Race-Conscious Intent.} A selective public magnet high school or state university eliminates racial checkboxes after \textit{Grutter}'s demise but adopts a facially race-neutral admissions preference for applicants from specific low-income ZIP codes and underfunded feeder schools, expressly hoping to maintain racial diversity. Under \textit{Grutter} (which encouraged experimenting with race-neutral alternatives, \href{https://supreme.justia.com/cases/federal/us/539/306/}{539 U.S. at 342}) and post-\textit{SFFA} doctrine, this sits on the razor's edge between permissible race-neutral socioeconomic targeting and impermissible intentional racial proxy discrimination.
\end{enumerate}

\section{Critique}
\textbf{Originalist, Colorblind, and Empirical Critiques:} Conservative and libertarian commentators criticized \textit{Grutter} for diluting strict scrutiny through unprecedented judicial deference to university administrators and for departing from the Fourteenth Amendment's individualist, colorblind command (\href{https://www.cato.org/supreme-court-review/2002-2003/principle-policy-public-university-admissions-grutter-v-bollinger}{Cato Supreme Court Review}). Following Justice Thomas's dissent, empirical scholar Richard H. Sander argued in \textit{A Systemic Analysis of Affirmative Action in American Law Schools}, 57 Stan. L. Rev. 367 (2004), that large admissions preferences produce academic ``mismatch''---placing minority students in schools where their entering credentials lag behind peers, allegedly leading to lower law school grades and higher bar-failure rates---a thesis vigorously contested by other empiricists (\href{https://scholarship.law.bu.edu/cgi/viewcontent.cgi?article=1325&context=faculty_scholarship}{Boston University Law Review Commentary}). Ultimately, in \href{https://supreme.justia.com/cases/federal/us/600/20-1199/}{\textit{Students for Fair Admissions, Inc. v. President and Fellows of Harvard College}}, 600 U.S. 181 (2023), the Supreme Court held that Harvard's and UNC's race-conscious admissions programs failed strict scrutiny because their diversity goals were unmeasurable, used race as a negative and a stereotype, and lacked a logical end point---effectively ending \textit{Grutter}'s era of race-conscious university admissions (\href{https://en.wikipedia.org/wiki/Grutter_v._Bollinger}{Wikipedia}).

\textbf{Progressive and Critical Race Critiques:} Progressive and critical race scholars also faulted \textit{Grutter} from the left. In his seminal 2003 essay \textit{Diversity's Distractions}, 103 Colum. L. Rev. 1622 (2003), Derrick Bell argued that grounding affirmative action in ``diversity'' rather than remedying structural racial subordination was a fragile ``distraction'' that centered elite institutions' interests (and benefits to white students) while diverting attention from poverty and K--12 inequities (\href{https://sclawreview.org/article/diversitys-distractions-revisited-the-case-of-latinx-in-higher-education/}{South Carolina Law Review}). Other scholars observed that \textit{Grutter} and \textit{Gratz} incentivized admissions obscurity---rewarding holistic opacity over transparent criteria (\href{https://texaslawreview.org/the-goose-and-the-gander-how-conservative-precedents-will-save-campus-affirmative-action/}{Texas Law Review}).

\textbf{Competing Constitutional Values:} \textit{Grutter} sits at the crossroads of \textit{anticlassification} (the principle that the Equal Protection Clause forbids the state from treating individuals differently based on race, regardless of motive) and \textit{antisubordination / institutional inclusion} (the principle that equal citizenship requires open pathways to elite education and civic leadership in a society still shaped by historical racial inequality).

\section{Key Quotations}
\begin{quote}
``Strict scrutiny is not `strict in theory, but fatal in fact.' ... Context matters when reviewing race-based governmental action under the Equal Protection Clause.''
\end{quote}
\noindent --- Justice Sandra Day O'Connor, Opinion of the Court, 539 U.S. at 326--327 (\href{https://supreme.justia.com/cases/federal/us/539/306/}{Justia}; \href{https://www.courtlistener.com/opinion/130155/grutter-v-bollinger/}{CourtListener}).

\begin{quote}
``In order to cultivate a set of leaders with legitimacy in the eyes of the citizenry, it is necessary that the path to leadership be visibly open to talented and qualified individuals of every race and ethnicity.''
\end{quote}
\noindent --- Justice Sandra Day O'Connor, Opinion of the Court, 539 U.S. at 332 (\href{https://supreme.justia.com/cases/federal/us/539/306/}{Justia}; \href{https://www.courtlistener.com/opinion/130155/grutter-v-bollinger/}{CourtListener}).

\begin{quote}
``We expect that 25 years from now, the use of racial preferences will no longer be necessary to further the interest approved today.''
\end{quote}
\noindent --- Justice Sandra Day O'Connor, Opinion of the Court, 539 U.S. at 343 (\href{https://supreme.justia.com/cases/federal/us/539/306/}{Justia}; \href{https://www.courtlistener.com/opinion/130155/grutter-v-bollinger/}{CourtListener}).

\begin{quote}
``Stripped of its `critical mass' veil, the Law School's program is revealed as a naked effort to achieve racial balancing.''
\end{quote}
\noindent --- Chief Justice William H. Rehnquist, dissenting, 539 U.S. at 379 (\href{https://supreme.justia.com/cases/federal/us/539/306/}{Justia}; \href{https://www.courtlistener.com/opinion/130155/grutter-v-bollinger/}{CourtListener}).

\begin{quote}
``The Constitution abhors classifications based on race, not only because those classifications can harm favored races or are based on illegitimate motives, but also because every time the government places citizens on racial registers and makes race relevant to the provision of burdens or benefits, it demeans us all.''
\end{quote}
\noindent --- Justice Clarence Thomas, concurring in part and dissenting in part, 539 U.S. at 353 (\href{https://supreme.justia.com/cases/federal/us/539/306/}{Justia}; \href{https://www.courtlistener.com/opinion/130155/grutter-v-bollinger/}{CourtListener}).

\end{document}
